%% kapa: exam=1 round=27 year=2016 period=1 form=A subject=economics subject_ko=경제학원론 seq=31 number=71
%% source: Q-Net 감정평가사 기출문제 (cst00309) artlSeq=5169402 file=extracted/1008/5169402/제27회 감정평가사 1차 문제/제27회 감정평가사 1교시문제(A형).pdf
%% conversion: kordoc 4.24.51 -> tools/gen1.py (auto); review=proofread-sonnet; flags=none
\begin{question}{71}
실물경기변동이론(real business cycle theory)에 관한 설명으로 옳은 것을 모두 고른 것은?
\begin{framedtext}
ㄱ. 임금 및 가격이 경직적이다. \par
ㄴ. 불경기에는 생산의 효율성이 달성되지 않는다. \par
ㄷ. 화폐의 중립성(neutrality of money)이 성립된다. \par
ㄹ. 경기변동은 시간에 따른 균형의 변화로 나타난다. \par
\end{framedtext}
\begin{choices}
  \item ㄱ, ㄴ
  \item ㄱ, ㄷ
  \item ㄴ, ㄷ
  \item ㄴ, ㄹ
  \item ㄷ, ㄹ
\end{choices}
\end{question}
\answer{5}
